\section{Security Testing}
\label{sec:security-testing}

El control A.14.3 de ISO 27001 responde una pregunta clave: \textit{¿Cómo sabemos que realmente es seguro?} Decir ``seguimos buenas prácticas'' no es evidencia. Hay que demostrarlo mediante pruebas \cite{iso27001}.

\subsection{Tipos de Pruebas de Seguridad}

\subsubsection{SAST --- Static Application Security Testing}

El análisis estático de seguridad analiza el código sin ejecutarlo:

\begin{center}
\small
\begin{tabular}{c}
\textbf{Código fuente} \\[2pt]
$\downarrow$ \\[2pt]
\textbf{Escáner} \\[2pt]
$\downarrow$ \\[2pt]
\textbf{Hallazgos}
\end{tabular}
\end{center}

\textbf{Qué busca:}
\begin{itemize}
    \item SQL Injection / NoSQL Injection.
    \item Hardcoded secrets (credenciales en código).
    \item Cross-Site Scripting (XSS).
    \item Uso inseguro de APIs.
\end{itemize}

\textbf{Herramientas para Workcodile-dev:}
\begin{itemize}
    \item \textbf{ESLint:} plugins \texttt{eslint-plugin-security}, \texttt{eslint-plugin-node}.
    \item \textbf{Semgrep:} análisis estático con reglas para Node.js.
    \item \textbf{SonarQube:} análisis completo de calidad y seguridad.
\end{itemize}

SAST permite detectar fallos temprano, antes de compilar o desplegar.

\subsubsection{DAST --- Dynamic Application Security Testing}

El análisis dinámico prueba la aplicación mientras está ejecutándose:

\begin{center}
\small
\begin{tabular}{c}
\textbf{Aplicación viva} \\[2pt]
$\downarrow$ \\[2pt]
\textbf{Scanner} \\[2pt]
$\downarrow$ \\[2pt]
\textbf{Ataques simulados}
\end{tabular}
\end{center}

\textbf{Qué busca:}
\begin{itemize}
    \item XSS reflejado y almacenado.
    \item SQL/NoSQL Injection.
    \item Headers de seguridad inseguros.
    \item Configuraciones débiles y CSRF.
\end{itemize}

\textbf{Herramientas para Workcodile-dev:}
\begin{itemize}
    \item \textbf{OWASP ZAP:} prueba automatizada de la app corriendo.
    \item \textbf{Burp Suite:} escaneo con proxy interceptador.
\end{itemize}

\subsubsection{Diferencia clave entre SAST y DAST}

\begin{itemize}
    \item \textbf{SAST ve el código fuente} --- análisis estático, sin ejecutar.
    \item \textbf{DAST ve el comportamiento} --- análisis dinámico, con la app corriendo.
    \item Ambos son complementarios.
\end{itemize}

\subsubsection{Code Review --- Revisión de Código}

La revisión humana es fundamental porque hay problemas que ninguna herramienta detecta.

\textbf{Ejemplo:}

\begin{lstlisting}[language=Java, caption={Lógica potencialmente incorrecta}]
if (user.role == "admin") {
  // Permitir acción
}
\end{lstlisting}

La lógica podría estar mal aunque el código sea sintácticamente correcto. Se debería usar \texttt{===} en lugar de \texttt{==}.

\textbf{En Workcodile-dev:}
\begin{itemize}
    \item Revisiones de código obligatorias antes de merge a \texttt{main}.
    \item Checklist de seguridad en cada Pull Request.
    \item Mínimo 1 aprobación antes de merge.
\end{itemize}

\subsubsection{Vulnerability Assessment}

Escaneo de bibliotecas, dependencias, servidores y contenedores:

\begin{itemize}
    \item \texttt{npm audit} --- dependencias con vulnerabilidades conocidas.
    \item \texttt{docker scan} --- vulnerabilidades en imágenes Docker.
    \item \texttt{trivy} --- escaneo completo de contenedores.
\end{itemize}

\subsubsection{Penetration Testing}

Un especialista intenta comprometer el sistema de forma controlada para demostrar impacto real.

\subsection{Acceptance Testing de Seguridad}

Es la parte más importante del control A.14.3. Antes de pasar a producción se definen criterios claros:

\begin{table}[h]
\centering
\small
\begin{tabular}{|l|l|}
\hline
\textbf{Criterio} & \textbf{Objetivo} \\
\hline
Vuln. críticas & 0 \\
Vuln. altas & 0 \\
Cobertura SAST & $\geq$ 80\% \\
Deps sin CVEs críticos & 100\% \\
Pruebas ejecutadas & Sí \\
Config. endurecida & Sí \\
\hline
\end{tabular}
\caption{Criterios de aceptación de seguridad}
\label{tab:criterios}
\end{table}

Si no cumple, no se despliega a producción.

\subsection{Pipeline CI/CD Seguro}

La materialización de los controles A.14.1, A.14.2 y A.14.3:

\begin{lstlisting}[caption={Pipeline CI/CD seguro}, basicstyle=\ttfamily\tiny, columns=fullflexible]
Git Push
  -> Build (npm install)
  -> Unit Tests (Jest / Vitest)
  -> SAST (ESLint + Semgrep)
  -> Dependency Scan (npm audit)
  -> DAST (OWASP ZAP)
  -> Security Gate
  -> Deploy Staging
  -> Acceptance Security Test
  -> Production
\end{lstlisting}

Este pipeline es la materialización práctica de:
\begin{itemize}
    \item \textbf{A.14.1} $\rightarrow$ proceso seguro de desarrollo.
    \item \textbf{A.14.2} $\rightarrow$ código seguro.
    \item \textbf{A.14.3} $\rightarrow$ validación de seguridad.
\end{itemize}
